\documentclass{article}

\usepackage{amsmath,amssymb}
\usepackage{xurl}
\usepackage[hidelinks]{hyperref}
\setlength{\emergencystretch}{2em}

% Shared commands for notes in docs/paper-gaps/.

\DeclareUnicodeCharacter{2080}{\ifmmode{}_{0}\else\textsubscript{0}\fi}
\DeclareUnicodeCharacter{2081}{\ifmmode{}_{1}\else\textsubscript{1}\fi}
\DeclareUnicodeCharacter{2082}{\ifmmode{}_{2}\else\textsubscript{2}\fi}
\DeclareUnicodeCharacter{2099}{\ifmmode{}_{n}\else\textsubscript{n}\fi}

\newcommand{\Lean}{\textsc{Lean}}
\ExplSyntaxOn
\NewDocumentCommand{\leanid}{m}{
  \ifmmode
    \textsf{\detokenize{#1}}
  \else
    \group_begin:
    \urlstyle{sf}
    \tl_set:Nn \l_tmpa_tl {#1}
    \tl_replace_all:Nnn \l_tmpa_tl {\_} {_}
    \exp_args:NV \path \l_tmpa_tl
    \group_end:
  \fi
}
\ExplSyntaxOff
\providecommand{\tr}{\operatorname{tr}}
\newcommand{\tnleanIssueBase}{https://github.com/LionSR/TNLean/issues/}
\newcommand{\tnleanPrBase}{https://github.com/LionSR/TNLean/pull/}
\newcommand{\tnleanDocBase}{https://sirui-lu.com/TNLean/docs/}
\newcommand{\ghissue}[1]{\href{\tnleanIssueBase#1}{GitHub issue \##1}}
\newcommand{\ghpr}[1]{\href{\tnleanPrBase#1}{PR \##1}}
\newcommand{\doclink}[1]{\href{\tnleanDocBase#1.html}{\path{#1}}}

% Dirac notation and the identity operator, as in blueprint/src/macros/common.tex.
% \providecommand keeps note-local definitions authoritative.
\usepackage{dsfont}
\providecommand{\ket}[1]{|#1\rangle}
\providecommand{\bra}[1]{\langle #1|}
\providecommand{\braket}[2]{\langle #1 | #2 \rangle}
\providecommand{\ketbra}[2]{|#1\rangle\!\langle #2|}
\providecommand{\Id}{\mathds{1}}
\providecommand{\one}{\mathds{1}}
\providecommand{\C}{\mathbb{C}}
\providecommand{\MN}[1]{M_{#1}(\C)}
\providecommand{\MD}{M_D(\C)}

% External referencing for paper sources.  The build helper writes sanitized
% auxiliary files under build/paper-aux/ and excludes duplicated source labels.
\usepackage{xr}

% Import a generated paper auxiliary file with a caller-supplied label prefix.
% The candidate paths support builds from docs/paper-gaps/, the repository root,
% and build/paper-aux/ itself.
\newcommand{\importpaperaux}[2]{%
  \IfFileExists{../../build/paper-aux/#2.aux}{%
    \externaldocument[#1]{../../build/paper-aux/#2}%
  }{%
    \IfFileExists{build/paper-aux/#2.aux}{%
      \externaldocument[#1]{build/paper-aux/#2}%
    }{%
      \IfFileExists{#2.aux}{%
        \externaldocument[#1]{#2}%
      }{}%
    }%
  }%
}

% General external reference macros
\newcommand{\paperref}[2]{\ref{#1-#2}}
\newcommand{\papereqref}[2]{(\ref{#1-#2})}

% CPSV16 source (1606.00608 / MPDO-22-12-17-2.tex) external document and shortcuts
\importpaperaux{cpsv16-}{1606.00608/MPDO-22-12-17-2-tnlean}
\newcommand{\cpsvref}[1]{\paperref{cpsv16}{#1}}
\newcommand{\cpsveqref}[1]{\papereqref{cpsv16}{#1}}

% Verdict marker: \gapnote{<kind>}{<status>}. Kind is the policy
% classification (clarification, local-correction, scope-restriction,
% unfaithful, false-source, open-gap); status is open, wip (elimination
% underway), resolved, or historical. Machine-read by the site index;
% prints a verdict line.
\newcommand{\gapnote}[2]{\par\noindent\textbf{Verdict:} #1 (#2).\par}


\title{Theorem Surface Audit for Canonical Form, BNT, and the Non-periodic MPS Fundamental Theorem}
\author{}
\date{2026-05-05}

\begin{document}
\maketitle
\gapnote{clarification}{open}

\section{Purpose}

This note compares the formal theorem surface with the canonical-form, basis of
normal tensors (BNT), and Fundamental Theorem statements in the source papers.
It distinguishes:
\begin{enumerate}
    \item statements named directly in the papers;
    \item mathematical results cited and used by the papers;
    \item auxiliary formal declarations that make implicit proof steps exact;
    \item stronger or alternative results that should remain outside the
          paper-level route unless that route requires them.
\end{enumerate}
The note classifies mathematical roles rather than tracking issues.

\section{Named source statements}

The canonical-form and BNT part of Ref.~\cite{Cirac2016MPDO_arXiv} contains a
small set of named source statements:
\begin{itemize}
    \item the definitions of normal tensors and canonical form
          \cite[lines~233--255]{Cirac2016MPDO_arXiv};
    \item the definition and characterization of a basis of normal tensors
          \cite[lines~271--280]{Cirac2016MPDO_arXiv};
    \item the definitions of injective and block-injective canonical form,
          together with the quantum Wielandt and
          block-injective-after-blocking assertions
          \cite[lines~317--345]{Cirac2016MPDO_arXiv};
    \item the proportional Fundamental Theorem and its equal-MPV corollary
          \cite[lines~347--356]{Cirac2016MPDO_arXiv}.
\end{itemize}
Section~IV of Ref.~\cite{Cirac2021Matrix} gives the same normal-tensor, BNT,
proportional, and equal-MPV results
\cite[lines~1828--1899]{Cirac2021Matrix}. These statements are the
paper-level targets and should remain prominent in the blueprint.

\section{Why the formalization has more declarations}

The source proof leaves several finite-dimensional and finite-index operations
implicit. A proof assistant requires separate declarations for those
operations, but their formal visibility does not make them additional
paper-level theorems.

\paragraph{Reduction statements.}
The source instruction to ``put the tensor in canonical form'' is decomposed
formally into invariant-subspace extraction, block-diagonal decomposition,
trace-preserving normalization, period-removing blocking, and normality of the
primitive blocks. These declarations expose the hypotheses and outputs needed
by later proofs. They do not add mathematical claims beyond the canonical-form
reduction in Refs.~\cite{Cirac2016MPDO_arXiv,Cirac2021Matrix}.

\paragraph{Common blocking statements.}
The source removes periodicity by blocking a common multiple of the periods
\cite{Cirac2016MPDO_arXiv}. The formal route must prove that the block length is
positive, that injectivity persists under further blocking, and that one block
length works simultaneously for a finite family of representative sectors.
These statements are auxiliary implementations of the source reduction.

\paragraph{Wielandt proof-chain statements.}
The proof of Theorem~1 in Ref.~\cite{Sanz2010Quantum} combines its Lemma~1,
eigenvalue extraction, vector spreading, and a fixed-length rank-one argument.
The source presents these ingredients as one proof, not as separate MPS
theorems. Formal conjunctions named ``Wielandt analysis'' or ``Wielandt
chain'' should therefore remain lemmas. The visible theorem should state the
cumulative-span or exact-length Wielandt conclusion and should distinguish the
stronger exact-length result from the cumulative-span bound.

\paragraph{BNT comparison statements.}
The BNT characterization in Ref.~\cite{Cirac2016MPDO_arXiv} chooses one
representative from each phase-gauge equivalence class and compares the power
sums of the associated weights. The formal proof separates this argument into
overlap decay, eventual linear independence, permutation rigidity,
sector-weight comparison, and phase absorption. A rectangular block-matching
theorem, in which the two sector decompositions may initially use different
bases, is a genuine formal result for the block-matching stage. The final
global weighted gauge equivalence remains the paper-level equal-MPV target.

The overlap-to-linear-independence and change-of-basis declarations support
permutation rigidity; they are not additional BNT theorems from the source.
For a tensor $S$, sector $j$, copy count $r_j^S$, and copy weights
$\mu_{j,q}^S$, define
\begin{align}
    c_{S,j}^{(N)}
        &=
        \sum_{q=1}^{r_j^S}
        (\mu_{j,q}^S)^N.
    \label{eq:tnlean_bnt_ft_theorem_surface_sector_coefficient}
\end{align}
For a second tensor $T$, the coefficient comparison yields the multiset
identity
\begin{align}
    \left\{\!\left\{\mu_{j,q}^S:1\leq q\leq r_j^S\right\}\!\right\}
        &=
        \left\{\!\left\{\mu_{j,q}^T:1\leq q\leq r_j^T\right\}\!\right\}.
    \label{eq:tnlean_bnt_ft_theorem_surface_weight_multiset_equality}
\end{align}
The double braces in
(\ref{eq:tnlean_bnt_ft_theorem_surface_weight_multiset_equality}) denote
multisets, so repeated weights retain their multiplicities. BNT linear
independence separates the coefficient arrays in
(\ref{eq:tnlean_bnt_ft_theorem_surface_sector_coefficient}), and the
Newton--Girard identities recover the weights. Formal declarations that merely
package these two algebraic steps in sector-data notation are support lemmas,
even if their names mention the Fundamental Theorem.

One former route first recovered the copy counts $r_j^S$ and $r_j^T$ by
extending eventual equality of positive power sums to exponent zero:
\begin{align}
    \sum_{q=1}^{r_j^S}
        (\mu_{j,q}^S)^0
        &=
        \sum_{q=1}^{r_j^T}
        (\mu_{j,q}^T)^0.
    \label{eq:tnlean_bnt_ft_theorem_surface_zero_power_identity}
\end{align}
It then applied the equal-cardinality Newton--Girard recovery statement. This
route is valid when every weight is nonzero: the geometric-sequence
extrapolation lemma derives
(\ref{eq:tnlean_bnt_ft_theorem_surface_zero_power_identity}) from eventual
equality at positive lengths. It is not, however, the argument displayed in
Lemma~\texttt{Lem:app\_simple}, lines~1155--1163, of
Ref.~\cite{Cirac2016MPDO_arXiv}. After matching the BNT representatives, the
source applies an unequal-cardinality power-sum lemma directly to the nonzero
canonical-form coefficients.

This classification follows
\path{docs/paper-gaps/policy.tex}. The detailed comparison appears in
\path{docs/paper-gaps/cpsv16_power_sum_alternative_route.tex}.

\section{External lemmas need explanations}

When an MPS source cites a result without proving it, the formal exposition
should explain the result where it enters the proof. A citation alone does not
supply the mathematical interface. The explanation should state:
\begin{enumerate}
    \item the exact mathematical conclusion used;
    \item the hypotheses in the current MPS notation, including any blocking
          or reindexing;
    \item a short proof idea or a precise statement that the proof is imported
          from the cited source;
    \item the formal declaration supplying the result, if it has been
          formalized.
\end{enumerate}

This rule applies to the following external inputs.
\begin{itemize}
    \item The quantum Wielandt theorem derives finite-length matrix-algebra
          span from normality or primitivity
          \cite{Sanz2010Quantum}. The canonical-form exposition should state
          whether it uses cumulative span, exact-length span, injectivity of
          $\Gamma_L$, or block injectivity of a direct sum.
    \item The block-injective canonical-form bound uses the earlier MPS
          injectivity and block-injectivity argument in
          Ref.~\cite{PerezGarcia2007MPSReps}. Its formal statement should
          identify the common blocked physical alphabet and the direct-sum
          matrix algebra being spanned. The direct-sum proof input is recorded
          in \path{docs/paper-gaps/pgvwc07_direct_sum_input.tex}.
    \item The Newton--Girard identities are elementary finite-sequence input,
          but Ref.~\cite{Cirac2016MPDO_arXiv} does not prove them. When they
          recover coefficients, multiplicities, or weight multisets, the
          blueprint should explain the finite power-sum recursion.
          Declarations that only restate this recursion in BNT notation should
          remain lemmas.
    \item Finite-dimensional Perron--Frobenius and peripheral-spectrum facts
          for quantum channels are distinct inputs. A formal theorem must state
          whether it assumes irreducibility, primitivity, trace preservation, a
          faithful fixed point, or a rank-one peripheral projection.
\end{itemize}

The same rule supplies a retirement test. An external-looking declaration
should not remain a headline theorem in the paper-level blueprint if it has no
source citation, no mathematical explanation, and no current proof path that
uses it.

\section{Dependency-chain classification}

The non-periodic Fundamental Theorem route requires a strict separation between
source targets, formal intermediate results, and stronger conclusions.

\subsection{Final target trace-back}

The source target is the Fundamental Theorem for canonical-form MPS and its
equal-MPV corollary, after reduction to normal form and removal of periodicity
by blocking. In the notation of lines~347--356 of
Ref.~\cite{Cirac2016MPDO_arXiv}, the proportional theorem compares two
canonical-form tensors whose MPV families agree up to a scalar sequence. It
concludes a permutation of the blocks, phases, and blockwise gauges. The
equal-MPV corollary sets the scalar sequence identically equal to $1$, but its
proof still uses the BNT coefficient comparison.

The formal route has three distinct layers:
\begin{enumerate}
    \item The declaration
          \leanid{MPSTensor.fundamentalTheorem_singleBlock} is the
          single-block injective gauge theorem. It is source-level only for one
          injective block; it does not establish the multi-block BNT
          comparison.
    \item The declarations
          \leanid{MPSTensor.ft_sector_bnt_equal_sector_data} and
          \leanid{MPSTensor.ft_sector_bnt_equal_global_gauge} are the formal
          results closest to the BNT part of the source proof. They use raw
          sector weights and derive sector matching, equality of copy counts,
          and matched copy-weight identities on the \texttt{SectorBNT}
          surface. The equal-MPV theorem on that surface applies after the
          canonical-form supplier establishes the required
          unit-modulus-copy normalization.
    \item The declaration
          \leanid{MPSTensor.fundamentalTheorem_equalMPV_sectorDecomposition_hetero_of_sectorMatching}
          is an auxiliary rectangular comparison theorem. Given a sector-basis
          matching, linear independence, and equality of the total MPV
          families, it proves equality of the corresponding sector-weight
          multisets up to the phases in the matching. It is not the paper-level
          equal-MPV corollary.
\end{enumerate}

The current proportional theorem is
\leanid{MPSTensor.fundamentalTheorem_proportional_canonicalForm}.%
\footnote{Formal declaration:
    \leanid{MPSTensor.fundamentalTheorem_proportional_canonicalForm}.}
It is complete on the \texttt{SectorBNT} surface. Every representative has a
positive copy count and every copy weight is nonzero, as required by the
nonzero-coefficient convention in
\path{docs/paper-gaps/cpsv16_bnt_uniqueness_zero_coefficient.tex}. From the
proportional MPV hypothesis, it derives a bijection of sectors, matched bond
dimensions, unit phases, and invertible block gauges. Separate global-gauge
declarations strengthen this sectorwise conclusion and should not be
identified with the source theorem.

The declaration
\leanid{MPSTensor.CPSVCanonicalFormData.fundamentalTheorem_equal_ambient_canonicalForm}
covers the equal-ambient case. It assumes equal ambient bond dimensions and
separate weight normalization for both canonical-form data sets. Exact
reconstruction and a coisometry-alignment argument then yield ambient gauge
equivalence. This theorem does not remove the heterogeneous
ambient-dimension obstruction in the printed corollary of
Ref.~\cite{Cirac2016MPDO_arXiv}.

\paragraph{Conditional after-blocking statements.}
The after-blocking comparison theorems apply only after one has supplied the
source canonical-form reduction, BNT selection, block-injective span, and BNT
matching inputs. The common-phase-cover and BNT-cover formulations are
conditional statements whose hypotheses encode those inputs. They must remain
visibly conditional until the paper assumptions are shown to imply them. In
particular, finite-length span and phase-cover assumptions are not additional
source hypotheses; they are formal names for proof obligations that still need
to be obtained from the source argument in
Ref.~\cite{Cirac2016MPDO_arXiv}.

\paragraph{Data structures and transport lemmas.}
Objects such as \leanid{SectorBasisOverlapSpanHypotheses} store formal
comparison data. They are not paper theorems. Declarations that convert one
such structure into another are auxiliary unless they directly state one of
the mathematical claims in Refs.~\cite{Cirac2016MPDO_arXiv,Cirac2021Matrix}.

\paragraph{External finite-dimensional inputs.}
The route also names finite-dimensional results that the source uses without
proof: Newton--Girard recovery, Perron--Frobenius structure, Burnside--Jacobson
algebra generation, and the quantum Wielandt bounds
\cite{Cirac2016MPDO_arXiv,Sanz2010Quantum}. These results may be substantial
theorems independently of the MPS argument. Within this proof chain, however,
they are cited inputs. Their mathematical statements should be explained, and
declarations that merely restate them in MPS notation should normally remain
lemma-level.

\section{Current demotions}

The following declarations are auxiliary and should not be read as independent
paper-level theorems.
\begin{itemize}
    \item The explicit-coefficient equal-MPV declaration
          \leanid{fundamentalTheorem_equalMPV_of_explicit_coefficients} is
          lemma-level. It assumes coefficient arrays and their nonzero limits,
          so it is not the full equal-MPV corollary in
          Refs.~\cite{Cirac2016MPDO_arXiv,Cirac2021Matrix}.
    \item The former proportional result on the old strict-representative
          surface has been removed from the active formal route. It assumed
          explicit coefficient arrays, their limits, and nonvanishing, whereas
          the source theorem derives the BNT coefficient comparison from the
          proportional MPV hypothesis
          \cite{Cirac2016MPDO_arXiv}. This was an incorrect theorem surface,
          not merely a theorem-versus-lemma presentation issue.
    \item The unused fixed-size BNT Vandermonde MPV declarations and the unused
          scalar Vandermonde separation lemma have been removed. The checked
          block-separation proof uses mixed-transfer peeling, while the
          equal-case BNT weight recovery uses finite power sums.
    \item The equal-span permutation-rigidity statement with asymptotically
          orthonormal overlaps is a lemma. It provides a useful span argument,
          but the main proportional Fundamental Theorem in Chapter~11 uses the
          proportional BNT permutation theorem.
    \item The older theorem that converted coefficient-array comparison data
          into sector comparison without bookkeeping for leftover zero blocks
          has been removed. The active after-blocking route retains this
          bookkeeping because its formal equality relation includes the empty
          word. Ref.~\cite{Cirac2016MPDO_arXiv} does not use this terminology;
          it states only that canonical forms may contain zero blocks through
          the inequality
          \begin{align}
              \sum_k D_k
                  &\leq
                  D.
              \label{eq:tnlean_bnt_ft_theorem_surface_canonical_block_dimension}
          \end{align}
          See lines~217--219 of Ref.~\cite{Cirac2016MPDO_arXiv}.
    \item The declaration
          \leanid{power_sum_eq_implies_multiset_eq} is a lemma-level
          restatement of the elementary Newton--Girard input used in the
          equal-MPV proof. The former declaration
          \leanid{perBlock_sameMPV_of_equalMPV_CFBNT}, which extracted
          blockwise equality from the common-block equal-MPV lemma, has also
          been removed from the active Lean tree. These names identify the
          retired one-copy route and should not be counted as additional
          source-level Fundamental Theorems.
    \item The implication
          \leanid{sameMPV₂_implies_proportionalMPV₂} converts equal MPV
          families into proportional MPV families with proportionality
          constant $1$. It is lemma-level because it is a logical conversion
          between hypotheses, not a separate theorem from
          Refs.~\cite{Cirac2016MPDO_arXiv,Cirac2021Matrix}. The subscript in
          the identifier is the Unicode subscript used in the Lean declaration
          name.
    \item The BNT sector-weight recovery declarations
          \leanid{weight_multiset_eq_of_sameMPV_bnt} and
          \leanid{exists_copies_eq_and_weight_multiset_eq_of_sameMPV_bnt}
          are lemma-level. They formalize the coefficient identity
          \begin{align}
              c_{S,j}^{(N)}
                  &=
                  c_{T,j}^{(N)}
              \label{eq:tnlean_bnt_ft_theorem_surface_matched_sector_coefficients}
          \end{align}
          and the elementary recovery of copy counts and weight multisets.
          They do not add BNT theorems beyond the basis-normal-tensor
          comparison in Ref.~\cite{Cirac2016MPDO_arXiv}.
    \item The phase-matched sector-comparison adapters
          \leanid{fundamentalTheorem_equalMPV_sectorDecomposition_hetero_of_phaseMatch}
          and
          \leanid{fundamentalTheorem_equalMPV_sectorDecomposition_hetero_of_matched_basis}
          are lemma-level. Their legacy names describe rectangular comparison
          adapters that transport the same coefficient and multiset formulas
          across a supplied phase or matched-basis identification. The
          theorem-level surface begins only after the BNT hypotheses produce
          the matching data.
\end{itemize}

\section{Statements to keep visibly conditional}

Several declarations remain useful but must retain explicit hypotheses because
they either begin from more data than the source theorem or prove a stronger
conclusion.
\begin{itemize}
    \item The conditional after-blocking Fundamental Theorem assumes the
          post-blocking sector hypotheses. It is not the unconditional
          canonical-form theorem in
          Refs.~\cite{Cirac2016MPDO_arXiv,Cirac2021Matrix}.
    \item The explicit-coefficient equal case assumes coefficient arrays
          instead of deriving the BNT power-sum data from equality of MPVs.
    \item The conditional proportional global-gauge theorem assumes
          scalar-free matched-sector coefficient identities. Under projective
          proportionality, the common length-dependent scalar $c_N$ need not
          be geometric. The global-gauge conclusion is therefore an
          intentionally stronger theorem, not unfinished work in the
          sectorwise proportional theorem.
\end{itemize}

\section{Retirement criterion}

An auxiliary declaration should be removed or moved outside the main blueprint
route when all of the following conditions hold:
\begin{enumerate}
    \item it is neither a statement in a cited paper nor a cited external
          result;
    \item no current proof of a paper-level target uses it;
    \item it assumes stronger finite-length, span, or synchronization data than
          the source proof requires.
\end{enumerate}
A declaration that remains useful only as a formal intermediate step should be
a lemma or a visibly conditional theorem. It should not replace the
paper-level Fundamental Theorem.

\bibliographystyle{alpha}
\bibliography{references}

\end{document}
